\documentclass[11pt]{article}
\usepackage[margin=1in]{geometry}
\usepackage{amsmath,amssymb,amsthm}
\usepackage{graphicx}
\usepackage{hyperref}

\newtheorem{theorem}{Theorem}
\newtheorem{remark}{Remark}

\newcommand{\R}{\mathbb R}
\newcommand{\E}{\mathbb E}
\newcommand{\cL}{\mathcal L}
\newcommand{\cH}{\mathcal H}

\title{A non-adapted subcase of the antisymmetric weak-subordination problem}
\author{arXiv automatic math run \texttt{fa\_banach\_001}}
\date{2026-06-18}

\begin{document}
\maketitle

\section*{Source and open target}

The source paper is Yaroslavtsev, \emph{Fourier multipliers and weak
differential subordination of martingales in UMD Banach spaces},
arXiv:1703.07817v5 \cite{Yaroslavtsev1703}.  In Section 5.2 it records
Conjecture 5.7, a linear-in-\(\beta_{p,X}\) weak differential subordination
estimate for continuous martingales.  The paper then proves Theorem 5.10 for
stochastic integrals \((\Phi A)\cdot W^H\) when \(A\) is self-adjoint, and
states the following open antisymmetric variant.

\begin{center}
\includegraphics[width=0.96\textwidth]{figures/open_problem_crop.png}
\end{center}

Thus the open point is the adapted-integrand estimate
\[
 \|((\Phi A)\cdot W^H)_\infty\|_{L^p(\Omega;X)}
 \leq C_p\beta_{p,X}\|A\|\|(\Phi\cdot W^H)_\infty\|_{L^p(\Omega;X)}
\]
for antisymmetric \(A\), especially the two-dimensional rotation which yields
the Hilbert-transform estimate.

\section*{Result}

The following partial result shows that the obstruction in Remark 5.12 is not
the antisymmetry of \(A\) by itself.  It is the interaction with Brownian-path
adaptivity of the integrand.

\begin{theorem}
Let \(H\) be a real Hilbert space, \(X\) a finite-dimensional real Banach
space, \(1<p<\infty\), and \(W^H\) an \(H\)-cylindrical Brownian motion.  Let
\(\mathcal G\) be a sigma-field independent of \(W^H\).  Let
\(\Phi:[0,T]\times\Omega\to \cL(H,X)\) be elementary and
\(\mathcal B([0,T])\otimes\mathcal G\)-measurable.  Then for every
\(A\in\cL(H)\),
\[
 \bigl(\E\|((\Phi A)\cdot W^H)_T\|^p\bigr)^{1/p}
 \leq
 \|A\|\,\bigl(\E\|(\Phi\cdot W^H)_T\|^p\bigr)^{1/p}.
\]
Consequently, for deterministic or initially random Brownian-independent
integrands, the antisymmetric case in Remark 5.12 holds with constant \(1\)
for every contraction \(A\).
\end{theorem}

\section*{Proof Intuition}

For deterministic \(\Phi\), both \((\Phi A)\cdot W^H\) and
\(\Phi\cdot W^H\) are centered Gaussian vectors in the finite-dimensional
Banach space \(X\).  If \(\|A\|\leq 1\), the covariance form of the first is
dominated by the covariance form of the second.  The missing covariance is
itself a positive covariance form, so one may add an independent Gaussian
remainder \(R\) to \((\Phi A)\cdot W^H\) and recover the law of
\(\Phi\cdot W^H\).  Jensen's inequality for the convex function
\(x\mapsto\|x\|^p\) then gives the estimate.  Conditioning on the initial
randomness gives the same argument for \(\mathcal G\)-measurable \(\Phi\).

The source problem remains hard precisely because a genuinely adapted
\(\Phi\) destroys this terminal Gaussian comparison.

\section*{Proof}

The case \(A=0\) is trivial.  By scaling \(A\), it is enough to treat
\(\|A\|\leq 1\).  First suppose that
\(\Phi\) is deterministic and elementary.  Set
\[
 M=(\Phi\cdot W^H)_T,\qquad N=((\Phi A)\cdot W^H)_T.
\]
Both are centered Gaussian \(X\)-valued random variables.  For \(x^*\in X^*\),
their covariance forms are
\[
 q_M(x^*)=\E|\langle M,x^*\rangle|^2
   =\int_0^T \|\Phi(t)^*x^*\|_H^2\,dt
\]
and
\[
 q_N(x^*)=\E|\langle N,x^*\rangle|^2
   =\int_0^T \|A^*\Phi(t)^*x^*\|_H^2\,dt.
\]
Since \(\|A\|\leq 1\),
\[
 q_N(x^*)\leq q_M(x^*)\qquad (x^*\in X^*).
\]
Thus \(q_R:=q_M-q_N\) is a positive semidefinite quadratic form on \(X^*\).
Because \(X\) is finite-dimensional, there exists a centered Gaussian
\(X\)-valued random variable \(R\), independent of \(N\), whose covariance
form is \(q_R\).  The centered Gaussian variables \(N+R\) and \(M\) have the
same covariance form, hence the same law.  Therefore, using Jensen's
inequality conditionally on \(N\),
\[
 \E\|N\|^p
 = \E\bigl\|\E(N+R\mid N)\bigr\|^p
 \leq \E\bigl(\E(\|N+R\|^p\mid N)\bigr)
 = \E\|N+R\|^p
 = \E\|M\|^p .
\]
This proves the assertion for deterministic elementary \(\Phi\) and
\(\|A\|\leq 1\).  Scaling gives the displayed estimate for arbitrary bounded
\(A\).

Now let \(\Phi\) be
\(\mathcal B([0,T])\otimes\mathcal G\)-measurable with \(\mathcal G\)
independent of \(W^H\).  Conditioning on \(\mathcal G\) freezes \(\Phi\) as a
deterministic elementary integrand, so the preceding estimate holds
conditionally.  Integrating the conditional inequality gives the theorem.

\section*{Verification notes and limitations}

The proof applies to the non-adapted, or initially randomized but
Brownian-independent, subcase.  It does not prove Conjecture 5.7 and does not
prove the adapted antisymmetric analogue of Theorem 5.10 from the source
paper.  In the adapted case, \(\Phi\) can depend on the Brownian path driving
the integral; the terminal integrals are no longer simply comparable centered
Gaussian vectors after conditioning, and the independent-remainder argument
has no direct extension.

The result is nevertheless a clean partial answer to Remark 5.12: for
deterministic integrands the estimate holds not only for antisymmetric \(A\),
but for every contraction \(A\), with constant \(1\) rather than
\(\beta_{p,X}\).

Local cheap-index search found no prior packet or attempt for arXiv:1703.07817.
A bounded web search on June 18, 2026 for arXiv:1703.07817, weak differential
subordination, the antisymmetric case, and the Hilbert-transform linear
estimate found the source paper, the follow-up martingale-decomposition paper
arXiv:1706.01731, and related strongly orthogonal martingale work
arXiv:1812.08049, but no existing run record for this deterministic subcase.

\begin{thebibliography}{9}

\bibitem{Yaroslavtsev1703}
I. Yaroslavtsev.
\newblock \emph{Fourier multipliers and weak differential subordination of
martingales in UMD Banach spaces}.
\newblock arXiv:1703.07817v5, 2017.

\end{thebibliography}

\end{document}
